\section{Lab 4: Gazebo Harmonic Simulation with Plugins and ROS 2 Bridge}

\subsection{Objective}
\begin{itemize}
	\item Explore the \texttt{<gazebo>} tags used to extend a URDF for Gazebo Harmonic.
	\item Add and understand a differential-drive plugin (and basic sensor plugins) in the robot description.
	\item Use the ROS~2--Gazebo bridge so ROS~2 nodes can talk to Gazebo.
	\item Teleoperate the robot from the keyboard and verify motion/outputs in Gazebo.
\end{itemize}

\subsection{Requirements}
\begin{itemize}
	\item \textbf{OS:} Ubuntu 22.04
	\item \textbf{ROS:} ROS 2 Humble
	\item \textbf{Simulator:} Gazebo Harmonic (gz sim 8)
	\item \textbf{Packages/Tools:} \texttt{ros\_gz\_sim}, \texttt{ros\_gz\_bridge}, \texttt{teleop\_twist\_keyboard}
\end{itemize}

\subsection{Theory}

\subsubsection{Gazebo plugins (what they are)}

Gazebo Harmonic lets us attach simulation-only settings to a robot description using \texttt{<gazebo>} tags. These tags do not change the real URDF structure; they only tell Gazebo how to simulate the model. Using them, we can:
\begin{itemize}
	\item Override visual/material properties for Gazebo
	\item Define friction/contact parameters
	\item Attach sensor plugins (cameras, LiDAR, IMU, etc.)
	\item Attach control plugins such as differential drive controller
\end{itemize}

\subsubsection{Differential drive controller (why we need it)}

For a differential drive robot, the plugin:
\begin{itemize}
	\item Subscribes to \texttt{cmd\_vel} topic (geometry\_msgs/Twist)
	\item Converts linear and angular velocity commands into left/right wheel velocities
	\item Uses wheel separation and wheel radius for calculations
	\item Publishes odometry data back to ROS 2
\end{itemize}

\subsubsection{ROS~2--Gazebo bridge (why bridging is needed)}

The \texttt{ros\_gz\_bridge} package creates a communication bridge between ROS 2 topics and Gazebo topics. This is necessary because:
\begin{itemize}
	\item Gazebo uses its own message format (Gazebo messages)
	\item ROS 2 uses standard ROS message types
	\item The bridge translates messages bidirectionally between these two systems
\end{itemize}

\subsection{Procedure}

\subsubsection{Prerequisites (installed tools)}

Before starting, we confirmed that the required tools were installed and that Gazebo could launch properly:
\begin{itemize}
	\item \texttt{ros\_gz\_sim} for spawning and running the robot in Gazebo
	\item \texttt{ros\_gz\_bridge} for bridging topics like \texttt{/clock}, \texttt{/cmd\_vel}, and \texttt{/odom}
	\item A teleop node (\texttt{teleop\_twist\_keyboard}) for publishing velocity commands
\end{itemize}

\subsubsection{Adding sensor plugins (overview)}

Next, we added basic sensor plugins (for example IMU/camera) using \texttt{<gazebo>} tags. The goal was not to build a full perception stack, but to understand how Gazebo plugins are declared and how their outputs appear as Gazebo topics (which we can bridge to ROS~2 when needed).

\subsubsection{Adding the differential drive plugin}

We added a differential drive controller plugin to enable motion control:

\begin{lstlisting}[language=XML]
<!-- Differential drive plugin (Gazebo Harmonic) -->
<gazebo>
  <plugin filename="gz-sim-diff-drive-system" 
          name="gz::sim::systems::DiffDrive">

    <!-- Left wheel joint name (must match the joint in URDF) -->
    <left_joint>left_wheel_joint</left_joint>

    <!-- Right wheel joint name (must match the joint in URDF) -->
    <right_joint>right_wheel_joint</right_joint>

    <!-- Distance between the wheels (meters) -->
    <wheel_separation>0.45</wheel_separation>

    <!-- Wheel radius (meters) -->
    <wheel_radius>0.1</wheel_radius>

    <!-- Odometry publishing frequency (Hz) -->
    <odom_publish_frequency>50</odom_publish_frequency>

    <!-- Topic where Twist commands are received -->
    <topic>cmd_vel</topic>
    
    <!-- Odometry topic -->
    <odom_topic>odom</odom_topic>

    <!-- Frame IDs for TF transforms -->
    <frame_id>odom</frame_id>
    <child_frame_id>base_link</child_frame_id>
  </plugin>
</gazebo>
\end{lstlisting}

\textbf{Plugin parameters (what each line means):}
\begin{itemize}
	\item \texttt{filename="gz-sim-diff-drive-system"}: Loads the built-in differential drive plugin
	\item \texttt{<left\_joint>} / \texttt{<right\_joint>}: Joint names that the plugin will control
	\item \texttt{<wheel\_separation>0.45}: Distance between wheel centers (left wheel at y=0.225, right at y=-0.225)
	\item \texttt{<wheel\_radius>0.1}: Radius of the wheels (cylinder radius from URDF)
	\item \texttt{<odom\_publish\_frequency>50}: Publishes odometry at 50 Hz
	\item \texttt{<topic>cmd\_vel</topic>}: Gazebo topic to receive velocity commands
	\item \texttt{<odom\_topic>odom</odom\_topic>}: Gazebo topic to publish odometry data
	\item \texttt{<frame\_id>odom}: Parent frame for odometry TF
	\item \texttt{<child\_frame\_id>base\_link}: Child frame for odometry TF
\end{itemize}

\textbf{Important:} The wheel separation and radius must match the physical robot geometry defined in the URDF. Incorrect values will cause the robot to move incorrectly or not move at all.

\subsubsection{Creating the Gazebo Launch File}

We created a launch file \texttt{gazebo\_sim.launch.py} to start Gazebo and spawn the robot. This launch file coordinates multiple nodes and processes to create a complete simulation environment.

\begin{lstlisting}[language=Python]
#!/usr/bin/env python3
import os
from launch import LaunchDescription
from launch.actions import ExecuteProcess, TimerAction
from launch_ros.actions import Node
from launch_ros.substitutions import FindPackageShare

def generate_launch_description():
    pkg_share = FindPackageShare('package_description').find('package_description')
    urdf_file = os.path.join(pkg_share, 'urdf', 'robot2_urdf.urdf')
    
    with open(urdf_file, 'r') as infp:
        robot_description = infp.read()
    
    # Launch Gazebo with empty world
    gazebo = ExecuteProcess(
        cmd=['gz', 'sim', 'empty.sdf', '-r'],
        output='screen'
    )
    
    # Spawn robot in Gazebo after delay
    spawn_robot = TimerAction(
        period=5.0,
        actions=[
            Node(
                package='ros_gz_sim',
                executable='create',
                arguments=[
                    '-topic', '/robot_description',
                    '-name', 'my_robot',
                    '-z', '0.5'
                ],
                output='screen'
            )
        ]
    )
    
    # Publish robot description
    robot_state_publisher = Node(
        package='robot_state_publisher',
        executable='robot_state_publisher',
        name='robot_state_publisher',
        output='screen',
        parameters=[{
            'robot_description': robot_description,
            'use_sim_time': True
        }]
    )
    
    # ROS-Gazebo Bridge
    bridge = Node(
        package='ros_gz_bridge',
        executable='parameter_bridge',
        arguments=[
            '/clock@rosgraph_msgs/msg/Clock[gz.msgs.Clock',
            '/cmd_vel@geometry_msgs/msg/Twist@gz.msgs.Twist',
            '/odom@nav_msgs/msg/Odometry@gz.msgs.Odometry'
        ],
        output='screen'
    )

    # Return all steps in order
    return LaunchDescription([
        gazebo,
        robot_state_publisher,
        spawn_robot,
        bridge
    ])
\end{lstlisting}

\subsubsection{Understanding the Gazebo Launch File Components}

\textbf{1. Package and URDF Setup:}
\begin{lstlisting}[language=Python]
pkg_share = FindPackageShare('package_description').find('package_description')
urdf_file = os.path.join(pkg_share, 'urdf', 'robot2_urdf.urdf')

with open(urdf_file, 'r') as infp:
    robot_description = infp.read()
\end{lstlisting}

\begin{itemize}
	\item \texttt{FindPackageShare}: Locates the installed package directory
	\item \texttt{os.path.join()}: Constructs the full path to the URDF file
	\item \texttt{with open()}: Reads the entire URDF file into a string variable
	\item This string will be passed to multiple nodes that need the robot description
\end{itemize}

\textbf{2. Launching Gazebo Simulator:}
\begin{lstlisting}[language=Python]
gazebo = ExecuteProcess(
    cmd=['gz', 'sim', 'empty.sdf', '-r'],
    output='screen'
)
\end{lstlisting}

\begin{itemize}
	\item \texttt{ExecuteProcess}: Runs Gazebo as a separate system process (not a ROS node)
	\item \texttt{'gz', 'sim'}: Command to launch Gazebo Harmonic
	\item \texttt{'empty.sdf'}: Loads an empty world file (no obstacles or terrain)
	\item \texttt{'-r'}: Auto-start flag - simulation begins running immediately
	\item \texttt{output='screen'}: Prints Gazebo logs to the terminal for debugging
\end{itemize}

\textbf{Why ExecuteProcess?} Gazebo is not a ROS node, it's a standalone simulator. We use \texttt{ExecuteProcess} to launch external applications, while \texttt{Node} is only for ROS nodes.

\textbf{3. Robot Spawning with Delay:}
\begin{lstlisting}[language=Python]
spawn_robot = TimerAction(
    period=5.0,
    actions=[
        Node(
            package='ros_gz_sim',
            executable='create',
            arguments=[
                '-topic', '/robot_description',
                '-name', 'my_robot',
                '-z', '0.5'
            ],
            output='screen'
        )
    ]
)
\end{lstlisting}

\begin{itemize}
	\item \texttt{TimerAction(period=5.0)}: Waits 5 seconds before executing
	\item \textbf{Why the delay?} Gazebo needs time to fully initialize before accepting spawn requests
	\item \texttt{ros\_gz\_sim/create}: Executable that inserts the robot model into Gazebo
	\item \texttt{'-topic', '/robot\_description'}: Reads URDF from this ROS topic
	\item \texttt{'-name', 'my\_robot'}: Assigns a unique name in Gazebo
	\item \texttt{'-z', '0.5'}: Spawns robot 0.5 meters above ground (prevents clipping into floor)
\end{itemize}

\textbf{4. Robot State Publisher:}
\begin{lstlisting}[language=Python]
robot_state_publisher = Node(
    package='robot_state_publisher',
    executable='robot_state_publisher',
    name='robot_state_publisher',
    output='screen',
    parameters=[{
        'robot_description': robot_description,
        'use_sim_time': True
    }]
)
\end{lstlisting}

\begin{itemize}
	\item Publishes the robot description to \texttt{/robot\_description} topic
	\item Publishes TF transforms showing where each link is in 3D space
	\item \texttt{'robot\_description': robot\_description}: Passes the URDF string as a parameter
	\item \texttt{'use\_sim\_time': True}: \textbf{Critical!} Uses Gazebo's simulation time instead of system clock
\end{itemize}

\textbf{Why use\_sim\_time is important:} In simulation, time can be paused, slowed down, or sped up. All nodes must use the same time source (Gazebo's clock) for synchronization. Without this, sensor timestamps and transforms would be mismatched.

\textbf{5. ROS-Gazebo Bridge:}
\begin{lstlisting}[language=Python]
bridge = Node(
    package='ros_gz_bridge',
    executable='parameter_bridge',
    arguments=[
        '/clock@rosgraph_msgs/msg/Clock[gz.msgs.Clock',
        '/cmd_vel@geometry_msgs/msg/Twist@gz.msgs.Twist',
        '/odom@nav_msgs/msg/Odometry@gz.msgs.Odometry'
    ],
    output='screen'
)
\end{lstlisting}

The bridge is the most critical component for ROS 2-Gazebo integration. Let's understand each bridge in detail:

\subsubsection{Understanding the ROS-Gazebo Bridge}

The bridge is critical for communication between ROS 2 and Gazebo. Each bridge argument follows this syntax:

\begin{verbatim}
/topic_name@ros2_msg_type@gazebo_msg_type    (bidirectional)
/topic_name@ros2_msg_type[gazebo_msg_type    (Gazebo to ROS)
/topic_name@ros2_msg_type]gazebo_msg_type    (ROS to Gazebo)
\end{verbatim}

\textbf{Why do we need bridges?}

Gazebo Harmonic and ROS 2 are completely separate systems:
\begin{itemize}
	\item \textbf{Gazebo} uses its own message format (Gazebo Transport / gz.msgs)
	\item \textbf{ROS 2} uses standard ROS message types (geometry\_msgs, nav\_msgs, etc.)
	\item They cannot directly communicate without translation
	\item The bridge acts as a translator, converting messages between these two formats
\end{itemize}

Without the bridge, ROS 2 nodes cannot send commands to Gazebo or receive sensor data from it.

\textbf{Bridge 1: Clock Synchronization}
\begin{lstlisting}[language=Python]
'/clock@rosgraph_msgs/msg/Clock[gz.msgs.Clock'
\end{lstlisting}

\begin{itemize}
	\item \textbf{Direction:} Gazebo → ROS (one-way, indicated by \texttt{[})
	\item \textbf{Purpose:} Synchronizes all ROS nodes with Gazebo's simulation time
	\item \textbf{ROS 2 message:} \texttt{rosgraph\_msgs/msg/Clock} - standard ROS clock message
	\item \textbf{Gazebo message:} \texttt{gz.msgs.Clock} - Gazebo's time representation
	\item \textbf{Why needed:} When simulation is paused/slowed/sped up, all nodes must use the same time
	\item \textbf{Published to:} \texttt{/clock} topic in ROS 2
	\item \textbf{Subscribers:} Any node with \texttt{use\_sim\_time: True} parameter
\end{itemize}

\textbf{What happens without clock bridge?}
\begin{itemize}
	\item Sensor data timestamps wouldn't match robot motion
	\item TF transforms would have incorrect timestamps
	\item Controllers would compute wrong velocities
	\item Navigation algorithms would fail
\end{itemize}

\textbf{Bridge 2: Velocity Commands}
\begin{lstlisting}[language=Python]
'/cmd_vel@geometry_msgs/msg/Twist@gz.msgs.Twist'
\end{lstlisting}

\begin{itemize}
	\item \textbf{Direction:} Bidirectional (indicated by \texttt{@})
	\item \textbf{Primary flow:} ROS 2 → Gazebo (sending commands)
	\item \textbf{Purpose:} Allows ROS 2 nodes to control the robot in Gazebo
	\item \textbf{ROS 2 message:} \texttt{geometry\_msgs/msg/Twist} - contains linear and angular velocities
	\item \textbf{Gazebo message:} \texttt{gz.msgs.Twist} - Gazebo's velocity representation
	\item \textbf{Topic:} \texttt{/cmd\_vel}
	\item \textbf{Publishers:} ROS 2 teleop nodes, navigation stack, autonomous controllers
	\item \textbf{Subscribers:} Differential drive plugin in Gazebo (from URDF)
\end{itemize}

\textbf{Data flow:}
\begin{enumerate}
	\item User presses keyboard keys in \texttt{teleop\_twist\_keyboard}
	\item Teleop node publishes \texttt{geometry\_msgs/Twist} to \texttt{/cmd\_vel} (ROS 2)
	\item Bridge converts to \texttt{gz.msgs.Twist}
	\item Message sent to Gazebo \texttt{/cmd\_vel} topic
	\item Differential drive plugin reads it and moves the wheels
\end{enumerate}

\textbf{Bridge 3: Odometry Data}
\begin{lstlisting}[language=Python]
'/odom@nav_msgs/msg/Odometry@gz.msgs.Odometry'
\end{lstlisting}

\begin{itemize}
	\item \textbf{Direction:} Bidirectional (indicated by \texttt{@})
	\item \textbf{Primary flow:} Gazebo → ROS 2 (receiving sensor data)
	\item \textbf{Purpose:} Provides robot position and velocity estimates to ROS 2
	\item \textbf{ROS 2 message:} \texttt{nav\_msgs/msg/Odometry} - position, orientation, velocities
	\item \textbf{Gazebo message:} \texttt{gz.msgs.Odometry} - Gazebo's odometry format
	\item \textbf{Topic:} \texttt{/odom}
	\item \textbf{Publishers:} Differential drive plugin in Gazebo
	\item \textbf{Subscribers:} Navigation nodes, localization algorithms, visualization tools
\end{itemize}

\textbf{Odometry data includes:}
\begin{itemize}
	\item \texttt{position}: (x, y, z) coordinates in the \texttt{odom} frame
	\item \texttt{orientation}: Rotation as a quaternion
	\item \texttt{linear velocity}: How fast the robot is moving
	\item \texttt{angular velocity}: How fast the robot is rotating
	\item \texttt{covariance}: Uncertainty in the estimates
\end{itemize}

\textbf{Why odometry is essential:}
\begin{itemize}
	\item Navigation algorithms need to know where the robot is
	\item Path planning requires current position
	\item Localization fuses odometry with other sensors
	\item Visualization tools display robot trail and position
\end{itemize}

\subsubsection{Bridge Syntax Summary}

\begin{tabular}{|l|l|l|}
\hline
\textbf{Symbol} & \textbf{Direction} & \textbf{Use Case} \\
\hline
\texttt{@} & Bidirectional & Commands and sensor data \\
\texttt{[} & Gazebo → ROS & Sensor data, clock, simulation info \\
\texttt{]} & ROS → Gazebo & Control commands, service calls \\
\hline
\end{tabular}

\vspace{0.3cm}

\textbf{Key principle:} Each bridge creates a translation layer for ONE topic. If you have 10 different sensors, you need 10 different bridge arguments (plus clock and control).

\subsubsection{Launch Sequence and Timing}

The order in which components are launched in the \texttt{LaunchDescription} is important:

\begin{lstlisting}[language=Python]
return LaunchDescription([
    gazebo,                  # 1. Start Gazebo first
    robot_state_publisher,   # 2. Publish URDF to /robot_description
    spawn_robot,             # 3. Wait 5s, then spawn robot
    bridge                   # 4. Start bridge for communication
])
\end{lstlisting}

\textbf{Why this specific order?}

\begin{enumerate}
	\item \textbf{Gazebo first:} The simulator must be running before we can spawn anything into it
	\item \textbf{Robot state publisher second:} Must publish URDF to \texttt{/robot\_description} before the spawn node tries to read it
	\item \textbf{Spawn robot (with 5s delay):} 
	\begin{itemize}
		\item Waits for Gazebo to fully initialize its physics engine
		\item Prevents spawn failures from attempting too early
		\item Gives time for all Gazebo plugins to load
	\end{itemize}
	\item \textbf{Bridge last:} Can start anytime, but placing it last ensures all topics are ready
\end{enumerate}

\textbf{What happens if order is wrong?}
\begin{itemize}
	\item Spawn before Gazebo → ERROR: No Gazebo server found
	\item Spawn before robot\_state\_publisher → ERROR: No URDF on \texttt{/robot\_description}
	\item No TimerAction delay → Robot may fail to spawn or appear corrupted
\end{itemize}

\subsubsection{Building and Running}

\textbf{Build the package:}
\begin{verbatim}
cd ~/ros2_ws
colcon build --packages-select package_description
source install/setup.bash
\end{verbatim}

\textbf{Launch Gazebo simulation:}
\begin{verbatim}
ros2 launch package_description gazebo_sim.launch.py
\end{verbatim}

\textbf{In a new terminal, teleoperate the robot:}
\begin{verbatim}
source install/setup.bash
ros2 run teleop_twist_keyboard teleop_twist_keyboard
\end{verbatim}

\subsubsection{Visualization proof}

The simulation visualization and the active ROS~2 topics during teleoperation are shown below.

\begin{figure}[h]
	\centering
	\begin{minipage}[t]{0.58\textwidth}
		\centering
		\IfFileExists{images/gazebo_view.png}{%
			\includegraphics[width=\textwidth]{images/gazebo_view.png}%
		}{%
			\fbox{\parbox{\textwidth}{\centering
				\textbf{Gazebo screenshot missing.}\\
				Place your file at \texttt{images/gazebo\_view.png}.
			}}%
		}
		\caption*{(a) Gazebo Harmonic view}
	\end{minipage}\hfill
	\begin{minipage}[t]{0.38\textwidth}
		\centering
		\IfFileExists{images/ros2_topic_list.png}{%
			\includegraphics[width=\textwidth]{images/ros2_topic_list.png}%
		}{%
			\fbox{\parbox{\textwidth}{\centering
				\textbf{Topic list screenshot missing.}\\
				Place your file at \texttt{images/ros2\_topic\_list.png}.
			}}%
		}
		\caption*{(b) ROS 2 topic list}
	\end{minipage}
	\caption{Gazebo Harmonic teleoperation and corresponding ROS~2 topics.}
	\label{fig:lab4-gazebo-view-topics}
\end{figure}

\subsection{Output}
\begin{itemize}
	\item Created a comprehensive launch file that coordinates Gazebo, robot spawning, and ROS 2 bridge
	\item Added differential drive plugin to URDF with correct wheel parameters (separation: 0.45m, radius: 0.1m)
	\item Integrated sensor plugins (IMU and other sensors) for robot perception
	\item Configured three critical bridges: clock synchronization, velocity commands, and odometry data
	\item Successfully launched Gazebo with empty world and spawned the differential drive robot
	\item Verified robot spawning with 5-second delay to ensure Gazebo initialization
	\item Confirmed \texttt{use\_sim\_time: True} for proper time synchronization across all nodes
	\item Successfully teleoperated the robot using \texttt{teleop\_twist\_keyboard}
	\item Verified odometry data publishing from Gazebo to ROS 2 through the bridge
	\item Observed smooth robot motion with accurate wheel velocities and turning behavior
\end{itemize}

\subsection{Inference}
\begin{itemize}
	\item The required ROS~2--Gazebo packages (\texttt{ros\_gz\_sim} and \texttt{ros\_gz\_bridge}) must be installed for simulation and topic bridging
	\item \texttt{ExecuteProcess} is used for launching external applications like Gazebo, not ROS nodes
	\item \texttt{Node} is used for ROS 2 nodes like robot\_state\_publisher and bridge
	\item \texttt{TimerAction} is essential to delay robot spawning until Gazebo is fully initialized (5 seconds is typical)
	\item Launch file order matters: Gazebo → robot\_state\_publisher → spawn\_robot → bridge
	\item \texttt{use\_sim\_time: True} ensures all nodes synchronize with Gazebo's simulation clock
	\item Without clock synchronization, sensor timestamps and transforms would be mismatched
	\item Gazebo plugins extend URDF with simulation-specific behaviors (sensors, controllers)
	\item The differential drive plugin requires accurate wheel geometry parameters (separation and radius)
	\item Wheel separation must match the actual distance between wheel joint origins in the URDF (0.45m in our case)
	\item Wheel radius must match the cylinder radius defined in the wheel link geometry (0.1m in our case)
	\item The ROS-Gazebo bridge acts as a message translator between ROS 2 and Gazebo formats
	\item Bridge directionality matters: \texttt{@} (bidirectional), \texttt{[} (Gazebo→ROS), \texttt{]} (ROS→Gazebo)
	\item Clock bridge (\texttt{[}) is unidirectional from Gazebo to ROS for time synchronization
	\item Velocity command bridge (\texttt{@}) is bidirectional but primarily ROS→Gazebo
	\item Odometry bridge (\texttt{@}) is bidirectional but primarily Gazebo→ROS
	\item Each topic needs its own bridge argument - no single bridge handles all topics
	\item Missing any bridge will break that specific communication channel
	\item Using \texttt{cmd\_vel} as a standard interface allows easy swapping of control methods
	\item Missing dependencies in package.xml will cause build or runtime errors
\end{itemize}

\subsection{Conclusions}
We successfully integrated Gazebo Harmonic simulation with ROS 2 by creating a comprehensive launch file that orchestrates multiple components in the correct sequence. The launch file demonstrated the importance of proper timing (TimerAction), the distinction between processes (Gazebo) and nodes (ROS 2 components), and the critical role of the ROS-Gazebo bridge in message translation. We learned that each bridge argument handles one specific topic, converting between ROS 2 and Gazebo message formats. The three bridges we implemented (clock, cmd\_vel, odom) enabled complete control and feedback for our differential drive robot. Understanding the launch file structure, bridge syntax, and component dependencies provides a solid foundation for building more complex robotic simulations with additional sensors and controllers.

ab5